\section{Geometric transport and local work}
\label{c:core:transport}

The first gravitational question is local: how does a prescribed geometry
change transport and energy exchange? A second question is dynamical: which
equations determine that geometry? Keeping these questions distinct lets a
finite model establish an exact balance law before it is coupled to a
gravitational constraint system. Throughout this section the metric
coefficients are supplied functions. The scalar field evolves on them.

\subsection{A finite scalar action on a diagonal metric}

Take a periodic chain with $N\geq3$ sites, spacing $h>0$, and indices modulo
$N$. The continuum metric motivating its coefficients is
\[
 ds^2=-N(t,x)^2dt^2+a(t,x)^2dx^2,
 \qquad N(t,x)>0,\quad a(t,x)>0.
\]
The repeated letter $N$ in this metric is the lapse, whereas the integer
$N$ counts lattice sites. In the finite formulas below only the coefficient
names $w$, $c$, and $v$ are needed:
\[
 w_j=\frac{a_j}{N_j}>0,\qquad
 c_{j+1/2}=\frac{N_{j+1/2}}{a_{j+1/2}}>0,\qquad
 v_j=N_ja_jm^2\geq0.
\]
Thus $N/a$ is the coordinate characteristic speed of the continuum
massless equation. Assume these finite coefficient arrays are continuously
differentiable in time. For a real field $\varphi_j(t)$ set
$\Delta_j=\varphi_{j+1}-\varphi_j$ and define
\begin{equation}
 L_h=\frac h2\sum_jw_j\dot\varphi_j^2
 -\frac1{2h}\sum_jc_{j+1/2}\Delta_j^2
 -\frac h2\sum_jv_j\varphi_j^2.
 \label{c:core:scalar-action}
\end{equation}
Variations with compact time support give
\begin{equation}
 w_j\ddot\varphi_j+\dot w_j\dot\varphi_j
 =\frac{c_{j+1/2}\Delta_j-c_{j-1/2}\Delta_{j-1}}{h^2}
 -v_j\varphi_j.
 \label{c:core:scalar-evolution}
\end{equation}
The neighboring edge terms in the action produce the two terms in the
spatial numerator. The time integration by parts differentiates
$w_j\dot\varphi_j$, including the $\dot w_j\dot\varphi_j$ term.

\Needspace{20\baselineskip}
\begin{proposition}[Local energy, flux, and geometric work]
\label{c:core:scalar-work}
For every twice differentiable solution of
\eqref{c:core:scalar-evolution}, define
\begin{align}
 e_j&=\frac h2w_j\dot\varphi_j^2+\frac h2v_j\varphi_j^2
 +\frac{c_{j+1/2}}{4h}\Delta_j^2
 +\frac{c_{j-1/2}}{4h}\Delta_{j-1}^2,\nonumber\\
 F_{j+1/2}&=-\frac{c_{j+1/2}}{2h}\Delta_j
             (\dot\varphi_j+\dot\varphi_{j+1}),\nonumber\\
 s_j&=-\frac h2\dot w_j\dot\varphi_j^2
 +\frac h2\dot v_j\varphi_j^2
 +\frac{\dot c_{j+1/2}}{4h}\Delta_j^2
 +\frac{\dot c_{j-1/2}}{4h}\Delta_{j-1}^2.
 \label{c:core:scalar-densities}
\end{align}
Then
\begin{equation}
 \dot e_j+F_{j+1/2}-F_{j-1/2}=s_j.
 \label{c:core:scalar-balance}
\end{equation}
In particular, $E_h=\sum_je_j$ satisfies
$\dot E_h=\sum_js_j=-\partial_tL_h$, where the partial derivative holds
$\varphi$ and $\dot\varphi$ fixed.
\end{proposition}

\begin{proof}
Put $B_j=c_{j+1/2}\Delta_j$. Using the equation of motion, the derivative
of the kinetic and mass terms at site $j$ is
\[
 \frac{\dot\varphi_j}{h}(B_j-B_{j-1})
 -\frac h2\dot w_j\dot\varphi_j^2
 +\frac h2\dot v_j\varphi_j^2.
\]
The right edge contribution differentiates to
\[
 \frac{\dot c_{j+1/2}\Delta_j^2}{4h}
 +\frac{B_j}{2h}(\dot\varphi_{j+1}-\dot\varphi_j),
\]
and the left edge has the analogous expression. Adding these terms
leaves
$B_j(\dot\varphi_j+\dot\varphi_{j+1})/(2h)$ on the right and its
negative counterpart on the left, plus $s_j$. These are exactly the flux
terms in \eqref{c:core:scalar-balance}. Periodic summation cancels every
interface flux. Each edge coefficient derivative appears twice with
weight $1/(4h)$, giving $-\partial_tL_h$.
\end{proof}

Every $e_j$ is nonnegative under the stated coefficient assumptions. An
interval of cells $r,\ldots,s$ obeys the exact coarse balance
\[
 \frac{d}{dt}\sum_{j=r}^se_j
 +F_{s+1/2}-F_{r-1/2}=\sum_{j=r}^ss_j.
\]
Thus coarse energy is obtained by adding fine cell energies and retaining
only the exposed boundary fluxes. This is the same cancellation mechanism
as projector transfer, applied here to a classical field energy. Its
physical units and constitutive coefficients come from the action.

\begin{example}[Time dependent expansion]
Let the lapse equal one, let $a(t)>0$ be spatially uniform, and let $m=0$.
Then $w=a$, $c=1/a$, and
\[
 \dot E_h=-\frac h2\dot a\sum_j\dot\varphi_j^2
 -\frac{\dot a}{2ha^2}\sum_j\Delta_j^2.
\]
An increasing $a$ transfers energy out of this coordinate-time field
Hamiltonian. The supplied background receives no evolution equation
from this calculation. A self-consistent matter--geometry model must
also evolve its gravitational variables.
\end{example}

For a smooth field sampled at sites and a smooth coefficient sampled at
edge midpoints, the spatial expression is second-order consistent with
$\partial_x(c\partial_x\varphi)$. Indeed, expanding at $x_j$ gives
\[
 c(x_j+h/2)[\varphi(x_j+h)-\varphi(x_j)]
 -c(x_j-h/2)[\varphi(x_j)-\varphi(x_j-h)]
 =h^2(c\varphi_x)_x(x_j)+O(h^4).
\]
Uniform bounds on the requisite fourth derivatives make the remainder
uniform on a compact spatial region. The exact finite balance holds
without any Taylor approximation.

\gravityfigure{01_local_energy_work}{A six-cell example of prescribed
metric work. Each direct cell-energy derivative agrees with incoming
flux minus outgoing flux plus the coefficient work term. The curves
illustrate the exact finite identity for a prescribed background.}
\FloatBarrier



\subsection{Conformal curvature as an exact boost holonomy}

Consider a twice continuously differentiable real function $\phi$ on an
oriented $(t,x)$ rectangle, and the coframe
\[
 e^0=e^\phi dt,\qquad e^1=e^\phi dx,
 \qquad ds^2=e^{2\phi}(-dt^2+dx^2).
\]
Write
\[
 K=\begin{pmatrix}0&1\\1&0\end{pmatrix},\qquad
 \omega=(\phi_xdt+\phi_tdx)K.
\]
The coframe satisfies the torsion-free Cartan equations with this
connection. Since every connection value is a multiple of the same
matrix, $\omega\wedge\omega=0$. Therefore
\begin{equation}
 \Omega=d\omega=(\phi_{tt}-\phi_{xx})K\,dt\wedge dx,
 \qquad R=2e^{-2\phi}(\phi_{tt}-\phi_{xx}).
 \label{c:core:conformal-curvature}
\end{equation}
This fixes the scalar-curvature sign convention used in this monograph.

We use the archive's plus-sign edge convention: for each oriented
boundary edge define $U_e=\exp(+\int_e\omega)$. This specifies the
signed transport observable used below. Multiplying in boundary order
yields an exact identity
\begin{equation}
 W_\square
 =\exp\left(K\oint_{\partial\square}(\phi_xdt+\phi_tdx)\right)
 =\exp\left(\frac K2\int_\square e^{2\phi}R\,dt\,dx\right).
 \label{c:core:conformal-holonomy}
\end{equation}
To see the second equality directly, integrate $\phi_{tt}$ along the
two horizontal endpoints and $\phi_{xx}$ along the two vertical
endpoints, with the boundary orientation from $dt\wedge dx$. This is
the rectangular fundamental theorem of calculus. No path-order
correction survives because the edge exponents commute.

If the integrated coefficient is $\chi$, then
\[
 e^{\chi K}=\cosh\chi\,I+\sinh\chi\,K.
\]
This preserves the Lorentz form $\operatorname{diag}(-1,1)$ and has
positive eigenvalues $e^{\chi}$ and $e^{-\chi}$. Its real principal
matrix logarithm is $\chi K$. These are Lorentz boosts. The positive
Hilbert norm used for unitary color transport in the gauge monograph
is a different structure.

\begin{example}[Constant curvature density]
For $\phi(t,x)=b(t^2-x^2)/2$, one has
\[
 \phi_{tt}-\phi_{xx}=2b,
 \qquad R=4b\,e^{-2\phi},
 \qquad W_\square=e^{2b\,\operatorname{area}(\square)K}.
\]
The curvature density $e^{2\phi}R/2$ is constant although $R$ varies
with position. Adding oriented subrectangle integrals reconstructs
the coarse rapidity exactly.
\end{example}

The plus-sign observable is assigned the endpoint action
\[
 G(t,x)=e^{\chi(t,x)K},\qquad
 U_{a\to b}'=G(b)U_{a\to b}G(a)^{-1}.
\]
It corresponds to the signed connection rule
$\omega_+'=\omega_++d\chi\,K$. Intermediate endpoint factors cancel,
so the loop observable is conjugated at its base point. For the
Cartan convention $D=d+\omega$, a coframe transformation $e\mapsto Ge$
instead gives $\omega'=\omega-d\chi\,K$ and parallel transport
$\exp(-\int_e\omega)$. Its closed-loop rapidity is the negative of the
rapidity of $W_\square$ defined here. The fixed-frame curvature identity
\eqref{c:core:conformal-holonomy} is unchanged. Declaring these signs
avoids mixing the archive's edge observable with a different Cartan
parallel-transport convention.



\section{Weak gravity as a sign and consistency benchmark}
\label{c:core:weak-field}

Weak gravity provides checks on attraction, variational signs, and
optical observables. These checks are performed within the weak-field
model. They do not furnish nonlinear gravitational evolution.

\subsection{Radial Poisson equation and a finite variational problem}

Let $G>0$, let $\rho(r)$ be a continuous spherical density, and impose
regularity $\lim_{r\downarrow0}r^2\Phi'(r)=0$. The radial Poisson
equation gives
\[
 (r^2\Phi')'=4\pi G r^2\rho,
 \qquad M_{\rm enc}(r)=4\pi\int_0^rs^2\rho(s)\,ds,
 \qquad \Phi'(r)=\frac{GM_{\rm enc}(r)}{r^2}.
\]
The acceleration is $a_r=-\Phi'$. A positive density therefore gives
an inward radial acceleration.

Putting $u(r)=r\Phi(r)$ removes the radial first derivative:
\[
 r^2\Phi'=ru'-u,\qquad
 (r^2\Phi')'=ru'',\qquad u''=4\pi G r\rho.
\]
For fixed endpoint values, this equation is the stationary condition of
\[
 S[u]=\int_0^{R_0}\left[\frac{(u')^2}{2G}+4\pi r\rho u\right]dr.
\]
On $r_i=ih$, $i=0,\ldots,M$, with $Mh=R_0$, its finite counterpart is
\begin{equation}
 S_h[u]=\sum_{i=0}^{M-1}\frac{(u_{i+1}-u_i)^2}{2Gh}
 +4\pi h\sum_{i=1}^{M-1}r_i\rho_i u_i.
 \label{c:core:poisson-action}
\end{equation}

\begin{proposition}[Unique finite radial solution]
\label{c:core:poisson-minimum}
With $u_0,u_M$ fixed, $S_h$ has a unique minimizer. Its interior values
satisfy
\[
 \frac{u_{i+1}-2u_i+u_{i-1}}{h^2}=4\pi G r_i\rho_i.
\]
\end{proposition}

\begin{proof}
Differentiation with respect to $u_i$ gives
$(2u_i-u_{i-1}-u_{i+1})/(Gh)+4\pi h r_i\rho_i$. Setting it to zero
gives the equation. For a variation $v$ with $v_0=v_M=0$, the Hessian
quadratic form is
\[
 \frac1{Gh}\sum_{i=0}^{M-1}(v_{i+1}-v_i)^2.
\]
Its vanishing forces $v$ to be constant and hence zero. The finite
Hessian is positive definite, so the quadratic action has one
stationary point and that point is its unique minimizer.
\end{proof}

The boundary value $u_0=0$ allows a regular finite potential at the
origin. The quotient $u/r$ is used only for $r>0$, with its continuous
limit determined separately at $r=0$.

For the normalized Gaussian density
\[
 \rho(r)=\frac{M_0}{(2\pi)^{3/2}\sigma^3}e^{-r^2/(2\sigma^2)},
 \qquad \sigma>0,
\]
direct differentiation verifies
\begin{align*}
 M_{\rm enc}(r)&=M_0\left[
 \operatorname{erf}\left(\frac r{\sqrt2\sigma}\right)
 -\sqrt{\frac2\pi}\frac r\sigma e^{-r^2/(2\sigma^2)}\right],\\
 \Phi(r)&=-\frac{GM_0}{r}
 \operatorname{erf}\left(\frac r{\sqrt2\sigma}\right).
\end{align*}
The potential tends to zero at infinity and has the finite central
value $-\sqrt{2/\pi}\,GM_0/\sigma$. This removes the point-source
singularity while preserving the attractive sign.

\subsection{Clocks and straight-ray optical integrals}

In the stipulated first-order isotropic weak-field metric,
\[
 ds^2=-\left(1+\frac{2\Phi}{c_0^2}\right)c_0^2dt^2
       +\left(1-\frac{2\Phi}{c_0^2}\right)d\mathbf x^2,
 \qquad \left|\frac{\Phi}{c_0^2}\right|\ll1,
\]
a static clock has lapse $\mathcal N=\sqrt{1+2\Phi/c_0^2}$.
The static-frequency ratio in this approximation is
\[
 \frac{\nu_{\rm obs}}{\nu_{\rm emit}}
 =\frac{\mathcal N_{\rm emit}}{\mathcal N_{\rm obs}}
 =1+\frac{\Phi_{\rm emit}-\Phi_{\rm obs}}{c_0^2}
 +O\left(\frac{\Phi^2}{c_0^4}\right).
\]
For $\Phi(x,b)=-GM_0/\sqrt{x^2+b^2}$, $b>0$, evaluating the
first-order straight-ray integrals over $[-L,L]$ gives
\begin{align*}
 \Delta t_{\rm Shapiro}
 &=-\frac2{c_0^3}\int_{-L}^L\Phi(x,b)\,dx
 =\frac{4GM_0}{c_0^3}\operatorname{asinh}\left(\frac Lb\right),\\
 \alpha(L,b)
 &=\frac2{c_0^2}\left|\int_{-L}^L\partial_b\Phi(x,b)\,dx\right|
 =\frac{4GM_0L}{c_0^2b\sqrt{L^2+b^2}}.
\end{align*}
Here $\alpha$ is the magnitude of the deflection toward the source.
The limit $L\to\infty$ gives $4GM_0/(c_0^2b)$. The displayed
equalities evaluate the weak-field integrals exactly. The approximation
to the physical ray already entered through the metric expansion and
the undeviated integration path.



\section{Gowdy evolution and reconstruction of the remaining constraint}
\label{c:core:gowdy}

The next step is an inhomogeneous nonlinear model with two gravitational
field channels. We take the archived Gowdy reduction as specified by its
action and constraint densities. The results below prove the internal
variational and reconstruction identities of that reduced system.

Let $P,Q$ be real smooth fields on $I\times\T$, with
$\theta\in\R/(2\pi\Z)$ and $\gamma(\tau)=e^{-2\tau}$. Define
\begin{equation}
 S_G=\frac12\int_I\int_0^{2\pi}
 \left[P_\tau^2-\gamma P_\theta^2
 +e^{2P}(Q_\tau^2-\gamma Q_\theta^2)\right]d\theta\,d\tau.
 \label{c:core:gowdy-action}
\end{equation}
Time variations vanish at the endpoints of $I$ and spatial variations
are periodic. Integration by parts gives the two Euler--Lagrange
expressions
\begin{align}
 \mathcal E_P&=P_{\tau\tau}-\gamma P_{\theta\theta}
 -e^{2P}(Q_\tau^2-\gamma Q_\theta^2),\nonumber\\
 \mathcal E_Q&=Q_{\tau\tau}-\gamma Q_{\theta\theta}
 +2(P_\tau Q_\tau-\gamma P_\theta Q_\theta).
 \label{c:core:gowdy-euler}
\end{align}
In particular, the $Q$ equation can first be written as
$\partial_\tau(e^{2P}Q_\tau)-\partial_\theta(e^{2P}\gamma Q_\theta)=0$.
The stationary equations are $\mathcal E_P=\mathcal E_Q=0$.

Define the densities
\begin{align}
 H_G&=P_\tau^2+\gamma P_\theta^2
 +e^{2P}(Q_\tau^2+\gamma Q_\theta^2),\nonumber\\
 M_G&=2(P_\tau P_\theta+e^{2P}Q_\tau Q_\theta).
 \label{c:core:gowdy-densities}
\end{align}
The remaining reconstruction equations are
$\lambda_\tau=H_G$ and $\lambda_\theta=M_G$.

\Needspace{14\baselineskip}
\begin{proposition}[Local integrability and global periodicity]
\label{c:core:gowdy-integrability}
For twice continuously differentiable $P,Q$,
\begin{equation}
 \partial_\tau M_G-\partial_\theta H_G
 =2P_\theta\mathcal E_P+2e^{2P}Q_\theta\mathcal E_Q.
 \label{c:core:gowdy-identity}
\end{equation}
For a solution on $I\times\T$, a single-valued periodic $\lambda$ with
these two derivatives exists if and only if
$\int_0^{2\pi}M_G\,d\theta=0$ at one, and hence every, time. It is
unique up to an additive constant.
\end{proposition}

\begin{proof}
Differentiate the two densities. The terms with $P_\tau
P_{\tau\theta}$ and $Q_\tau Q_{\tau\theta}$ cancel, leaving
\[
\begin{aligned}
 \partial_\tau M_G-\partial_\theta H_G
 ={}&2P_\theta(P_{\tau\tau}-\gamma P_{\theta\theta})\\
 &+2e^{2P}\big[
 Q_\theta(Q_{\tau\tau}-\gamma Q_{\theta\theta})
 +2P_\tau Q_\tau Q_\theta
 -P_\theta Q_\tau^2-\gamma P_\theta Q_\theta^2\big].
\end{aligned}
\]
Substitute \eqref{c:core:gowdy-euler} and combine the two occurrences of
$\gamma P_\theta Q_\theta^2$. This gives
\eqref{c:core:gowdy-identity}. On a solution the identity vanishes,
so periodicity gives $d\int M_G\,d\theta/d\tau=0$.

If a periodic $\lambda$ exists, integrating $\lambda_\theta=M_G$
forces the zero-period condition. Conversely, fix $\tau_0$ and define
on the lifted strip
\[
 \lambda(\tau,\theta)=\lambda_0
 +\int_{\tau_0}^{\tau}H_G(s,0)\,ds
 +\int_0^{\theta}M_G(\tau,\eta)\,d\eta.
\]
Its $\theta$ derivative is $M_G$. Its $\tau$ derivative is $H_G$ by
the integrability identity and the fundamental theorem of calculus.
The zero-period condition makes it periodic. The difference of two
such functions has both derivatives zero, so is constant.
\end{proof}

This separates a local compatibility condition from a global topology
condition. Satisfying the evolution equations makes the reconstruction
one-form closed. A vanishing spatial period makes it the derivative of
a periodic potential.

\subsection{An exact polarized mode}

Set $Q=0$ and take an integer $n\geq1$. Define the entire function
\[
 J_0(x)=\sum_{m=0}^{\infty}\frac{(-1)^m(x^2/4)^m}{(m!)^2}.
\]
The series and its differentiated series converge uniformly on compact
sets. Coefficient comparison gives
$x^2J_0''+xJ_0'+x^2J_0=0$. Thus
\begin{equation}
 P(\tau,\theta)=A J_0(ne^{-\tau})\cos(n\theta)
 \label{c:core:gowdy-mode}
\end{equation}
is an exact polarized solution. Indeed, for $x=ne^{-\tau}$,
$\partial_\tau^2J_0(x)=x^2J_0''(x)+xJ_0'(x)=-x^2J_0(x)$,
whereas $P_{\theta\theta}=-n^2P$. Also
$M_G=2P_\tau P_\theta$ has zero spatial integral, so periodic
$\lambda$ can be reconstructed.

The Taylor channels can be separated without changing the physical
mode. Put $y=x^2/4$, $\mathcal B(y)=J_0(2\sqrt y)$, and for an integer
$r\geq2$ define
\[
 \mathcal B_{r,j}(y)=\sum_{q=0}^{\infty}
 \frac{(-1)^{rq+j}y^{rq+j}}{((rq+j)!)^2},
 \qquad j=0,\ldots,r-1.
\]
Then $\mathcal B=\sum_j\mathcal B_{r,j}$. With
$\vartheta=y\,d/dy$, the coefficient identity
\[
 \vartheta^2\mathcal B_{r,j}+y\mathcal B_{r,j-1}=0,
 \qquad j\pmod r,
\]
follows from $m^2/(m!)^2=1/((m-1)!)^2$ for $m\geq1$.
Consequently the separated Taylor channels are coupled. Their sum
satisfies the Bessel equation. This decomposition resolves series
order and does not introduce additional independent gravitational
degrees of freedom.



\section{An exact finite Gowdy constraint projection}
\label{c:core:gowdy-projection}

Let $h=2\pi/N$ and use the periodic centered derivative
\[
 (D_hf)_j=\frac{f_{j+1}-f_{j-1}}{2h},\qquad
 (\Delta_hf)_j=\frac{f_{j+1}-2f_j+f_{j-1}}{h^2}.
\]
Here $D_h^{\mathsf T}=-D_h$. With $U=\dot P$, $V=\dot Q$, the raw
semidiscrete vector field $f_h(\tau,z)$ for $z=(P,Q,U,V)$ is
\begin{align}
 \dot P&=U,&\dot Q&=V,\nonumber\\
 \dot U&=\gamma\Delta_hP+e^{2P}[V^2-\gamma(D_hQ)^2],\nonumber\\
 \dot V&=\gamma\Delta_hQ-2[UV-\gamma(D_hP)(D_hQ)].
 \label{c:core:gowdy-raw}
\end{align}
Products and exponentials in these formulas are componentwise.
The finite global momentum is
\begin{equation}
 C_h(z)=2h\sum_j[U_j(D_hP)_j+e^{2P_j}V_j(D_hQ)_j].
 \label{c:core:gowdy-momentum}
\end{equation}
It is linear in the velocity variables when the configuration $P,Q$
is fixed.

For the ordinary Euclidean gradients, skew-adjointness of $D_h$ gives
\begin{align*}
 \nabla_PC_h&=-2hD_hU+4h e^{2P}V(D_hQ),\\
 \nabla_QC_h&=-2hD_h(e^{2P}V),\\
 g_U:=\nabla_UC_h&=2hD_hP,
 &g_V:=\nabla_VC_h&=2h e^{2P}D_hQ.
\end{align*}
Put $G_h=(0,0,g_U,g_V)$ and $d_h=\norm{g_U}_2^2+\norm{g_V}_2^2$.
On the open region $d_h>0$, define
\[
 \Lambda_h=\frac{\nabla C_h\cdot f_h}{d_h},\qquad
 \widehat f_h=f_h-\Lambda_hG_h.
\]
The chain rule proves the exact tangency identity
\[
 \nabla C_h\cdot\widehat f_h
 =\nabla C_h\cdot f_h-\Lambda_hd_h=0.
\]
Every differentiable solution that remains in this region therefore
keeps $C_h$ constant.

\begin{proposition}[Exact endpoint momentum]
\label{c:core:gowdy-endpoint}
Start from $z_n$ and let a provisional step produce
$\widetilde z=(\widetilde P,\widetilde Q,\widetilde U,\widetilde V)$.
Form $\widetilde g_U,\widetilde g_V$ at its configuration and assume
$\widetilde d_h>0$. The correction
\[
 \mu=\frac{C_h(\widetilde z)-C_h(z_n)}{\widetilde d_h},\qquad
 z_{n+1}=(\widetilde P,\widetilde Q,
 \widetilde U-\mu\widetilde g_U,
 \widetilde V-\mu\widetilde g_V)
\]
satisfies $C_h(z_{n+1})=C_h(z_n)$ exactly. It is the correction of
smallest Euclidean velocity norm at the fixed configuration.
\end{proposition}

\begin{proof}
Linearity in the velocities gives
$C_h(z_{n+1})=C_h(\widetilde z)-\mu\widetilde d_h$.
For minimality, write $g=(\widetilde g_U,\widetilde g_V)$ and
$b=C_h(z_n)-C_h(\widetilde z)$. Every admissible velocity displacement
$v$ satisfies $g\cdot v=b$. Cauchy--Schwarz implies
$\norm v\geq|b|/\norm g$, with equality only at
$v=bg/\norm g^2$. This is the displayed correction.
\end{proof}

\Needspace{11\baselineskip}
\begin{proposition}[Fixed-mesh second-order time accuracy]
\label{c:core:gowdy-time-order}
Fix $h$. Suppose a smooth exact solution of
$\dot z=\widehat f_h(\tau,z)$ and its numerical steps remain in a
compact regular neighborhood where $d_h\geq d_0>0$. Assume the mixed time and state
derivatives of $\widehat f_h$ through order three and the state
derivatives of $C_h$ through order four are bounded there. Apply the endpoint correction
of Proposition~\ref{c:core:gowdy-endpoint} after explicit midpoint.
Then its local error is $O(\Delta\tau^3)$ and its global error on a
fixed time interval is $O(\Delta\tau^2)$.
\end{proposition}

\begin{proof}
The explicit midpoint expansion agrees with the Taylor expansion of
the exact flow through order two. Its provisional endpoint differs
from the exact endpoint by $O(\Delta\tau^3)$. Since the exact flow
keeps $C_h$ constant and $C_h$ has bounded first derivatives, the
provisional constraint defect is also $O(\Delta\tau^3)$. The lower
bound on the denominator and the upper bound on $g$ make the
correction $O(\Delta\tau^3)$. Thus the corrected local error has the
same order. Under the stated derivative bounds the corrected step
map has Lipschitz constant at most $1+L\Delta\tau$ for small steps.
If $e_n$ is its error against the exact solution, then
\[
 e_{n+1}\leq(1+L\Delta\tau)e_n+M\Delta\tau^3.
\]
Iterating this inequality for $n\Delta\tau\leq T$ bounds the error by
$MT e^{LT}\Delta\tau^2$ when $e_0=0$.
\end{proof}

These are regular-chart results. For a spatially homogeneous
configuration, $g_U=g_V=0$, and the displayed multiplier is
undefined even though $C_h=0$. Constants in the accuracy theorem
can depend on $h$. A claim about simultaneous spatial and temporal
refinement additionally needs bounds uniform in $h$.

The local compatibility diagnostic is a separate quantity. Define
$H_{G,h},M_{G,h}$ by replacing spatial derivatives in
\eqref{c:core:gowdy-densities} with $D_h$. If
$\mathcal E_{P,h}=\dot U-\dot U_{\rm raw}$ and
$\mathcal E_{Q,h}=\dot V-\dot V_{\rm raw}$, then
\[
 \mathcal I_h=\partial_\tau M_{G,h}-D_hH_{G,h}
 -2(D_hP)\mathcal E_{P,h}
 -2e^{2P}(D_hQ)\mathcal E_{Q,h}
\]
subtracts the explicit constraint-force contribution. Exact global
momentum preservation does not force $\mathcal I_h$ to vanish at each
site. Its refinement behavior tests whether the modified evolution
approaches the local continuum reconstruction identity. The archived
second-order numerical observations are experimental evidence for
that behavior, with their tested range recorded in the companion.



\section{Finite linearized ADM gravity}
\label{c:core:adm}

The symmetry-reduced model can now be compared with a finite
three-dimensional linearized constraint system. Its exact algebraic
closure rests on commuting linear derivatives. Nonlinear products
will introduce an additional requirement in the next section.

\subsection{Finite calculus and the quadratic action}

Let $V_h$ be a finite real inner-product space. Let
$\delta_1,\delta_2,\delta_3$ be commuting skew-adjoint operators on it:
\[
 \delta_i\delta_j=\delta_j\delta_i,
 \qquad\inner{\delta_iX}{Y}_h=-\inner{X}{\delta_iY}_h,
 \qquad\delta^2=\sum_i\delta_i^2.
\]
The periodic centered differences in three coordinate directions
provide one example. All tensor inner products below sum over both
indices, using the Frobenius inner product on symmetric tensors.

Let $h_{ij}=h_{ji}$ and $\pi_{ij}=\pi_{ji}$ take values in $V_h$.
Repeated spatial indices are summed. Put $h_{\rm tr}=h_{kk}$,
$\pi_{\rm tr}=\pi_{kk}$ and define
\begin{align}
 C&=\delta_i\delta_jh_{ij}-\delta^2h_{\rm tr},
 &C_i&=-2\delta_j\pi_{ij},\nonumber\\
 \mathcal G_{ij}(h)&=\frac12\big[
 \delta_k\delta_ih_{jk}+\delta_k\delta_jh_{ik}
 -\delta^2h_{ij}-\delta_i\delta_jh_{\rm tr}
 -\delta_{ij}C\big].
 \label{c:core:adm-operators}
\end{align}

\begin{lemma}[Linearized finite Bianchi and adjoint identities]
\label{c:core:adm-bianchi}
The spatial operator obeys $\delta_j\mathcal G_{ij}(h)=0$ and is
self-adjoint on symmetric tensors.
\end{lemma}

\begin{proof}
Twice its divergence expands as
\begin{align*}
 2\delta_j\mathcal G_{ij}
 ={}&\delta_i\delta_j\delta_kh_{jk}
 +\delta^2\delta_kh_{ik}-\delta^2\delta_jh_{ij}
 -\delta_i\delta^2h_{\rm tr}\\
 &-\delta_i\delta_k\delta_lh_{kl}
 +\delta_i\delta^2h_{\rm tr}.
\end{align*}
The terms cancel in three pairs. For self-adjointness, write
$B_i(a)=\delta_ja_{ij}$ and $A(a)=\delta_i\delta_ja_{ij}$. Two
applications of the adjoint identity give
\begin{align*}
 2\inner{a}{\mathcal G(b)}_h
 ={}&-2\inner{B(a)}{B(b)}_h
 +\sum_k\inner{\delta_ka}{\delta_kb}_h\\
 &-\inner{A(a)}{b_{\rm tr}}_h
 -\inner{a_{\rm tr}}{A(b)}_h
 -\sum_k\inner{\delta_ka_{\rm tr}}{\delta_kb_{\rm tr}}_h.
\end{align*}
This expression is symmetric in $a,b$.
\end{proof}

Set
\[
 T_h=\inner{\pi}{\pi}_h-\tfrac12\inner{\pi_{\rm tr}}{\pi_{\rm tr}}_h,
 \qquad U_h=\tfrac12\inner{h}{\mathcal G(h)}_h.
\]
With lapse perturbation $n$ and shift perturbation $n_i$, the action is
\begin{equation}
 S_h^{(2)}=\int\left[
 \inner{\pi}{\dot h}_h-T_h-U_h
 +\inner{n}{C}_h-\inner{n_i}{C_i}_h\right]dt.
 \label{c:core:adm-action}
\end{equation}
Variation of $n,n_i$ imposes $C=C_i=0$. The remaining equations are
\begin{align}
 \dot h_{ij}&=2\pi_{ij}-\delta_{ij}\pi_{\rm tr}
             +\delta_in_j+\delta_jn_i,\nonumber\\
 \dot\pi_{ij}&=-\mathcal G_{ij}(h)
              +\delta_i\delta_jn-\delta_{ij}\delta^2n.
 \label{c:core:adm-evolution}
\end{align}
The kinetic gradient is $2\pi_{ij}-\delta_{ij}\pi_{\rm tr}$.
Lemma~\ref{c:core:adm-bianchi} makes the potential gradient
$\mathcal G(h)$. The lapse gradient follows by two spatial
integrations by parts. Integrating the shift term once gives
$-\langle\delta_in_j+\delta_jn_i,\pi_{ij}\rangle_h$ in the
Lagrangian, fixing the sign in the first equation.

\Needspace{16\baselineskip}
\subsection{Constraint propagation and the canonical brackets}

\Needspace{11\baselineskip}
\begin{theorem}[Exact finite linearized closure]
\label{c:core:adm-closure}
For every solution of \eqref{c:core:adm-evolution},
\[
 \dot C=-\delta_iC_i,\qquad\dot C_i=0.
\]
Thus initially vanishing constraints remain zero. With the canonical
Frobenius bracket and smearings
$C[n]=\langle n,C\rangle_h$, $D[\xi]=\langle\xi_i,C_i\rangle_h$,
all three brackets $\{C[n],C[m]\}$, $\{D[\xi],D[\eta]\}$, and
$\{C[n],D[\xi]\}$ vanish identically.
\end{theorem}

\begin{proof}
The trace of the momentum part of $\dot h$ is $-\pi_{\rm tr}$.
Its contribution to $\dot C$ is therefore
$2\delta_i\delta_j\pi_{ij}=-\delta_iC_i$. The shift contribution is
\[
 \delta_i\delta_j(\delta_in_j+\delta_jn_i)
 -\delta^2(2\delta_kn_k)=0.
\]
For $\dot C_i$, the divergence of $\mathcal G$ vanishes by
Lemma~\ref{c:core:adm-bianchi}, and
$\delta_j(\delta_i\delta_jn-\delta_{ij}\delta^2n)=0$.
This proves the propagation equations and their stated implication.

For the bracket, use
$\{F,G\}=\langle\nabla_hF,\nabla_\pi G\rangle_h
-\langle\nabla_\pi F,\nabla_hG\rangle_h$.
Two scalar constraints depend only on $h$ and two momentum
constraints depend only on $\pi$, so their brackets vanish.
The gradients for the mixed bracket are
\[
 (\nabla_hC[n])_{ij}=\delta_i\delta_jn-\delta_{ij}\delta^2n,
 \qquad(\nabla_\pi D[\xi])_{ij}=\delta_i\xi_j+\delta_j\xi_i.
\]
Their contraction vanishes after one integration by parts because
the first tensor is divergence-free.
\end{proof}

The momentum generator changes the metric perturbation by
$h_{ij}\mapsto h_{ij}+\delta_i\xi_j+\delta_j\xi_i$.
The corresponding change of the scalar constraint is
$2\delta^2\delta_k\xi_k-2\delta^2\delta_k\xi_k=0$.
No product rule is used in these cancellations.

For a partition of fine spatial sites into coarse regions $A$, the
nonnegative diagnostic
\[
 E_A^{\rm con}=\sum_{x\in A}h^3
 \left(|C_x|^2+\sum_i|C_{i,x}|^2\right)
\]
adds exactly to the full fine diagnostic. This is an exact partition
identity. It does not assert that differentiating a coarse field
reconstructs its fine constraint values.



\subsection{Tensor waves, positive reduced energy, and resolved dispersion}

Choose $n=n_i=0$ and impose the transverse and traceless conditions
\[
 \delta_jh_{ij}=0,\quad h_{\rm tr}=0,
 \qquad\delta_j\pi_{ij}=0,\quad\pi_{\rm tr}=0.
\]
They are preserved by the evolution because
\[
 \mathcal G_{ij}(h)=-\tfrac12\delta^2h_{ij},\qquad
 \dot h=2\pi,\qquad\dot\pi=\tfrac12\delta^2h,
 \qquad\ddot h=\delta^2h.
\]
On this reduced sector the energy is
\[
 H_{\rm TT}=\norm{\pi}_h^2
 +\frac14\sum_k\norm{\delta_kh}_h^2\geq0.
\]
Its derivative is zero: substituting the two first-order equations
reduces it to
$\langle\pi,\delta^2h\rangle_h
+\sum_k\langle\delta_kh,\delta_k\pi\rangle_h=0$.

For a Fourier mode on the periodic centered lattice, let
$\widetilde k_i=\sin(k_ih)/h$. If $\widetilde{\mathbf k}\neq0$,
define
\[
 \Pi_{ij}=\delta_{ij}
 -\frac{\widetilde k_i\widetilde k_j}{|\widetilde{\mathbf k}|^2},
 \qquad
 \Lambda^{\rm TT}_{ij,kl}
 =\tfrac12(\Pi_{ik}\Pi_{jl}+\Pi_{il}\Pi_{jk}-\Pi_{ij}\Pi_{kl}).
\]
For a symmetric tensor $X$, its image is
$\Pi X\Pi-\Pi\operatorname{tr}(\Pi X)/2$. It is transverse and
traceless because $\Pi^2=\Pi$ and $\operatorname{tr}\Pi=2$.
Applying the map again leaves it unchanged. On the transverse
two-plane, symmetric tensors have dimension three and the trace
condition removes one dimension. Thus the projector has rank two.

The resulting frequency obeys
\begin{equation}
 \omega_h^2=\sum_i\frac{\sin^2(k_ih)}{h^2}.
 \label{c:core:adm-dispersion}
\end{equation}
For fixed $\mathbf k\neq0$ as $h\to0$,
\[
 \omega_h^2=|\mathbf k|^2-\frac{h^2}{3}\sum_i k_i^4+O(h^4),
 \qquad
 \omega_h=|\mathbf k|
 -\frac{h^2}{6|\mathbf k|}\sum_i k_i^4+O(h^4).
\]
These follow by expanding the sine and then the positive square root.
The regime keeps $\mathbf k$ fixed. It is not uniform up to the
Nyquist boundary.

The spatial operator in \eqref{c:core:adm-dispersion} is the sum of
the squared centered first derivatives. It differs from the
nearest-neighbor second-difference operator used for the raw Gowdy
equations. For even site counts, an alternating sequence has centered
derivative zero. Such checkerboard modes, as well as constant modes,
belong to the derivative kernel. At a zero symbol the above TT
projector is undefined and the rank-two argument does not apply.
Exact constraint propagation consequently coexists with a finite
spatial kernel that a physical continuum comparison must address.




\gravityfigure{04_adm_dispersion}{The centered derivative has a Nyquist
zero, while a fixed resolved mode approaches its continuum frequency
with second-order error. The left panel is the one-dimensional symbol;
the right panel tests a fixed mode under refinement.}

\section{The nonlinear pointwise obstruction and its algebraic alternative}
\label{c:core:nonlinear-obstruction}

The linearized cancellations use commuting derivatives but never ask
them to differentiate a product. A nonlinear diffeomorphism generator
does. The following scalar model isolates that issue without relying
on the more involved gravitational tensor constraints.

Let a periodic real lattice have canonical brackets
$\{q_a,p_b\}_h=\delta_{ab}/h$ and a skew matrix $D$. Define
\[
 \mathcal D_h[\xi]=h\sum_a\xi_ap_a(Dq)_a,
 \qquad
 \mathfrak L_D(f,g)=D(fg)-(Df)g-f(Dg).
\]
All products on this lattice are pointwise.

\begin{proposition}[Exact nonlinear defect identity]
\label{c:core:nonlinear-defect}
For arbitrary finite arrays $q,p,\xi,\eta$,
\begin{align}
 \{\mathcal D_h[\xi],\mathcal D_h[\eta]\}_h
 ={}&\mathcal D_h[\xi D\eta-\eta D\xi]\nonumber\\
 &+h\sum_ap_a\big[
 \xi_a\mathfrak L_D(\eta,Dq)_a
 -\eta_a\mathfrak L_D(\xi,Dq)_a\big].
 \label{c:core:nonlinear-defect-identity}
\end{align}
\end{proposition}

\begin{proof}
Coordinate differentiation gives
$\partial_p\mathcal D_h[\xi]=h\xi Dq$ and
$\partial_q\mathcal D_h[\xi]=-hD(\xi p)$.
Insert these in the canonical bracket and use $D^{\mathsf T}=-D$:
\[
 \{\mathcal D_h[\xi],\mathcal D_h[\eta]\}_h
 =h\sum_ap_a\left[\xi_aD(\eta Dq)_a
                  -\eta_aD(\xi Dq)_a\right].
\]
Expanding each derivative using the definition of $\mathfrak L_D$
cancels the two $\xi\eta D^2q$ terms. The remaining product-rule
terms give the first term of
\eqref{c:core:nonlinear-defect-identity}, and the two defect terms
give the second.
\end{proof}

\begin{theorem}[Finite pointwise derivation obstruction]
\label{c:core:pointwise-no-go}
A linear derivation on the complete pointwise algebra $\C^N$ is zero.
Thus a nonzero finite pointwise derivative cannot obey the undeformed
Leibniz rule on every pair of arrays.
\end{theorem}

\begin{proof}
Let $e_a$ be the coordinate idempotents. If $D$ is a derivation, then
$D(e_a)=D(e_a^2)=2e_aD(e_a)$. At coordinate $a$ this forces the
component to be twice itself, hence zero. At every other coordinate
the right side is zero. Therefore $D(e_a)=0$ for every $a$. Linearity
and the coordinate basis imply $D=0$.
\end{proof}

The theorem rules out a universal ordinary derivation on the full
finite pointwise algebra. For a particular generator, one must still
inspect the contracted defect in
\eqref{c:core:nonlinear-defect-identity}. A nonzero Leibniz defect
does not by itself imply that every contraction is nonzero.

\begin{example}[An exact bracket failure]
\label{c:core:bracket-negative-control}
Use $N=5$, $h=1$, and $(Df)_a=(f_{a+1}-f_{a-1})/2$.
Take $q=e_0$, $p=\xi=e_1$, and $\eta=e_2$. Then
$(Dq)_1=-1/2$, $(D\eta)_1=1/2$, and $\eta Dq=0$.
The direct bracket in the preceding proof is zero, while
\[
 \mathcal D_h[\xi D\eta-\eta D\xi]=-\tfrac14.
\]
The defect is therefore exactly $+1/4$. This finite witness shows
that the naive centered scalar constraint does fail to close
universally in the proposed undeformed form.
\end{example}

\subsection{Smooth resolved arrays and the matrix corner route}

For samples of sufficiently smooth periodic functions, the centered
derivative obeys
\[
 D_hf=f'+\frac{h^2}{6}f'''+O(h^4).
\]
If $f,g$ have uniformly bounded derivatives through order five, its
Leibniz defect consequently satisfies
\begin{equation}
 \mathfrak L_{D_h}(f,g)
 =\frac{h^2}{2}(f''g'+f'g'')+O(h^4).
 \label{c:core:smooth-defect}
\end{equation}
For fixed smooth $p,\xi,\eta,q$, with sufficiently many additional
bounded derivatives of $q$ to control $D_hq$, the contraction in
\eqref{c:core:nonlinear-defect-identity} is $O(h^2)$ on a fixed periodic
interval. Each summand is uniformly $O(h^2)$ and $hN$ is the interval
length. The coordinate-idempotent witness above is outside this
uniformly smooth regime.

A spectral derivative has an exact product rule on trigonometric
polynomials when every product remains in the represented band without
aliasing. Repeated products require further headroom. These restricted
domains differ from the complete pointwise algebra. The later resolved
constraint theorem includes explicit domain bounds.

For the full matrix algebra, choose a Hermitian matrix $K$ and set
$\delta_KX=i[K,X]$. Direct multiplication gives
\begin{equation}
 \delta_K(XY)=(\delta_KX)Y+X(\delta_KY).
 \label{c:core:inner-derivation}
\end{equation}
It also obeys $\delta_K(X^*)=(\delta_KX)^*$,
$\operatorname{Tr}\delta_KX=0$, and is skew-adjoint for the
Hilbert--Schmidt inner product. The last identity follows by moving
$K$ through the trace. These properties supply an exact finite
product calculus on a noncommutative algebra.


